\chapter{Introdução}
\section{Contextualização}

Com evoluções tecnológicas constantes, o uso de técnicas de inteligência artificial (IA), principalmente algoritmos de aprendizado de máquina (AM), tem se tornado cada vez mais comum em diversas áreas. Atualmente, as principais técnicas de IA utilizadas são baseadas em redes neurais e transformadores, que são capazes de aprender padrões complexos a partir de grandes volumes de dados. No entanto, essas técnicas muitas vezes resultam em modelos que são considerados "caixas-pretas", ou seja, cuja lógica interna é difícil de interpretar e compreender totalmente. Em algumas áreas, como na medicina, a interpretabilidade dos modelos é crucial, pois permite um entendimento mais profundo dos fenômenos capturados pelos dados.

Dessa forma, outras técnicas de IA, como a regressão simbólica, tem ganhado destaque como forma de obter modelos mais interpretáveis. A regressão simbólica é uma técnica que busca encontrar expressões matemáticas que melhor se ajustam aos dados observados e que respeitem um conjunto pré-determinado de operações matemáticas. Algoritmos de regressão simbólica são capazes de gerar modelos matemáticos que podem ser facilmente interpretados e analisados, o que os torna uma alternativa interessante para aplicações onde a transparência é importante. Várias técnicas podem ser utilizadas para implementar a regressão simbólica, incluindo algoritmos evolutivos ou técnicas híbridas que combinam redes neurais com algoritmos evolutivos.

\section{Objetivos}
O objetivo deste trabalho é fazer uma revisão da literatura sobre regressão simbólica, apresentando suas principais características, técnicas e aplicações. Além disso, será desenvolvida uma interface web para facilitar a exploração de algoritmos de regressão simbólica, permitindo que usuários possam experimentar com diferentes conjuntos de dados e parâmetros do algoritmo.

\section{Justificativa}
Atualmente, ferramentas como TuringBot e HeuristicLab podem ser utilizadas para criação de modelos de regressão simbólica. No entanto, essas ferramentas possuem funcionalidades limitadas e devem ser executadas localmente, o que pode dificultar o uso em computadores com recursos limitados. A interface web proposta neste trabalho visa superar essas limitações, oferecendo uma plataforma acessível, na qual o treinamento dos modelos é feito remotamente, permitindo que usuários possam explorar a regressão simbólica de forma mais intuitiva e acessível.

Além disso, essa plataforma permitirá que usuários explorem soluções por meio de uma linguagem de consulta, similar ao SQL, para selecionar padrões matemáticos de interesse dentro do espaço de soluções proposto pelo algoritmo, garantindo que conhecimentos prévios sobre os fenômenos estudados possam ser incorporados ao processo de modelagem.

